\section{Closure architecture B: Yukawas, masses, flavor, confinement, and strong CP}

The second closure-architecture block is the transport compiler. Its role is to turn the carrier packet and the family/occupancy closure into actual masses, mixings, hadronic admissibility, and strong-CP control. The logic is again staged: first the discrete phase operator and the transport kernel, then path-length masses and mixing data, then confinement and strong CP as transport-side admissibility statements.

\subsection{Discrete phase operator and transport grammar}

Let
\[
\omega=e^{2\pi i/3},
\qquad
U_6=\mathrm{diag}(1,\omega,\omega^2,-1,-\omega,-\omega^2).
\]
This makes the phase grammar concrete rather than rhetorical. The six admissible phases are shown as a finite hexagonal transport geometry in \Cref{fig:z6-circle}: the cusp values live on the same real axis as the positive-real phase, the admissible gap is visible directly, and the holonomy phase can be read as a triangle on the same discrete transport base.

\begin{remark}[Hexagonal spectral geometry]
The operator $U_6$ is a unitary circulant with eigenphases $e^{in\pi/3}$ for $n=0,\dots,5$ and is unitarily equivalent to the one-step translation operator on a six-site ring. The transport kernel $D_y=y\,\mathbf 1-\delta_{\mathrm{ph}}U_6$ is therefore the resolvent denominator of a finite quantum walk, equivalently a tight-binding transfer problem, on the discrete $C_6$ cycle. The regularized inverse introduced below plays the role of the corresponding Green kernel. CKM and PMNS mixings, mass hierarchies, and the Möbius quotients of \Cref{eq:mobius-quotient} are thus not loose analytic fits but spectral-theory statements on a regular hexagon. The near-pole hierarchy at $y\approx\delta_{\mathrm{ph}}$ becomes the resonance of a hexagonal resolvent rather than tuned parameter choice.
\end{remark}

\begin{figure}[t]
\centering
\scriptsize
\begin{tikzpicture}[scale=1.0,>=Latex]
\coordinate (O) at (0,0);
\foreach \name/\ang in {p0/0,p1/60,p2/120,p3/180,p4/240,p5/300} {
    \coordinate (\name) at ({2*cos(\ang)},{2*sin(\ang)});
}
\draw[thick] (p0) -- (p1) -- (p2) -- (p3) -- (p4) -- (p5) -- cycle;
\draw[gray!35, dashed] (O) -- (p0);
\draw[gray!35, dashed] (O) -- (p1);
\draw[gray!35, dashed] (O) -- (p2);
\draw[gray!35, dashed] (O) -- (p3);
\draw[gray!35, dashed] (O) -- (p4);
\draw[gray!35, dashed] (O) -- (p5);
\foreach \P/\lab in {p0/{1},p1/{e^{i\pi/3}},p2/{e^{2i\pi/3}},p3/{-1},p4/{e^{-2i\pi/3}},p5/{e^{-i\pi/3}}} {
    \filldraw[blue!60!black] (\P) circle (2.3pt);
}
\node[anchor=west] at (2.22,0.02) {$1$};
\node[anchor=west] at (1.2,1.95) {$e^{i\pi/3}$};
\node[anchor=east] at (-1.2,1.95) {$e^{2i\pi/3}$};
\node[anchor=east] at (-2.2,0.02) {$-1$};
\node[anchor=east] at (-1.2,-1.95) {$e^{-2i\pi/3}$};
\node[anchor=west] at (1.2,-1.95) {$e^{-i\pi/3}$};

\draw[->, thick, red!70!black] (1.0,1.2) arc[start angle=48,end angle=-48,radius=1.35];
\node[red!70!black] at (2.7,0.3) {CP: $\theta\mapsto-\theta$};

\filldraw[orange!80!black] (0.67,0) circle (2.3pt);
\filldraw[orange!80!black] (1.33,0) circle (2.3pt);
\filldraw[green!50!black] (1.1998,0) circle (2.6pt);
\node[below] at (0.67,-0.02) {$\frac13$};
\node[below] at (1.33,-0.02) {$\frac23$};
\node[above] at (1.1998,0.02) {$\delta_{\mathrm{ph}}$};
\draw[<->, thick] (1.1998,-0.28) -- (1.33,-0.28);
\node[below] at (1.2649,-0.3) {$\epsilon_f$};
\node[align=left, anchor=west] at (2.6,-0.85) {near pole at\\$y\approx\delta_{\mathrm{ph}}$\\largest hierarchy gain};

\begin{scope}[shift={(4.35,-1.1)}]
    \coordinate (A) at (0,1.0);
    \coordinate (B) at (-0.95,-0.55);
    \coordinate (C) at (0.95,-0.55);
    \draw[thick, orange!75!black] (A) -- (B) -- (C) -- cycle;
    \filldraw[orange!75!black] (A) circle (1.7pt) (B) circle (1.7pt) (C) circle (1.7pt);
    \node[above] at (A) {$1$};
    \node[left] at (B) {$2$};
    \node[right] at (C) {$3$};
    \node[align=center] at (0,-1.15) {$\delta_{\mathrm{CKM}}=\arg\tr\Pi_\triangle$\\holonomy triangle};
\end{scope}
\end{tikzpicture}
\caption{Finite transport geometry on the hexagon $C_6$. The six $U_6$ phases form a discrete transport base, the admissible cusps $y\in\{1,\frac13,\frac23\}$ sit on the positive-real axis together with $\delta_{\mathrm{ph}}$, the smallest gap is $\epsilon_f=|\frac23-\delta_{\mathrm{ph}}|$, and the same transport grammar also carries the holonomy triangle phase.}
\label{fig:z6-circle}
\end{figure}

\subsection{Regularized Yukawa kernels and path-length masses}

\begin{definition}[Canonical transport saddle]
The transport pole $\delta_{\mathrm{ph}}$ is defined as the unique stationary point in the retained interval $(1/3,2/3)$ of the determinant transport functional
\[
\mathcal F_{\mathrm{tr}}(\delta)
:=
\sum_{y\in\{1,1/3,2/3\}}
\Bigl[-\log\det\!\bigl(D_{y,\delta}^\dagger D_{y,\delta}+\varepsilon^2\bigr)\Bigr]
+
\frac{36}{(\phiz)^2}\left(\delta-\frac35\right)^2.
\]
Here
\[
D_{y,\delta}:=y\,\mathbf 1-\delta U_6.
\]
Taking the regularization to $\varepsilon\to0$ and solving the exact stationarity condition on the retained branch gives
\begin{equation}
\label{eq:transport-bridge-offset}
\partial_\delta \mathcal F_{\mathrm{tr}}(\delta_{\mathrm{ph}})=0,
\qquad
\delta_{\mathrm{ph}}\in\left(\frac13,\frac23\right),
\qquad
\delta_{\mathrm{ph}}\approx \num{0.599882}.
\end{equation}
\end{definition}

Here $\delta_{\mathrm{ph}}\approx \num{0.599882}$ lies slightly below $3/5$, so $\epsilon_f=|\frac23-\delta_{\mathrm{ph}}|\approx \num{0.06678}$. Both are treated as retained benchmark-side transport numbers, not as independent hard-kernel constants. The full determinant-level variational derivation is deferred to the supplementary layer, while the regularized transport operator used here is
\begin{equation}
\label{eq:yukawa-kernel}
\begin{aligned}
D_y&=y\,\mathbf 1-\delta_{\mathrm{ph}}U_6,\\
\mathcal Y_y^{(\varepsilon)}&=(D_y^\dagger D_y+\varepsilon^2)^{-1},
\end{aligned}
\end{equation}
and the singular limit $\varepsilon\to0$ is taken only after admissibility has been restricted to the gap-stable sector defined below.

\begin{proposition}[Finite hexagonal resolvent]
\label{prop:finite-hexagonal-resolvent}
Whenever
\[
y^6\neq \delta_{\mathrm{ph}}^6,
\]
the transport denominator has the exact inverse
\begin{equation}
\label{eq:finite-hexagonal-resolvent}
D_y^{-1}
=
\frac{1}{y^6-\delta_{\mathrm{ph}}^6}
\sum_{m=0}^{5}y^{5-m}\delta_{\mathrm{ph}}^{\,m}U_6^m.
\end{equation}
Thus the $C_6$ resolvent is finite and algebraically closed rather than an infinite formal Green-series expansion.
\end{proposition}

\begin{proof}
Using $U_6^6=\mathbf 1$,
\[
\bigl(y\,\mathbf 1-\delta_{\mathrm{ph}}U_6\bigr)
\sum_{m=0}^{5}y^{5-m}\delta_{\mathrm{ph}}^{\,m}U_6^m
=
y^6\mathbf 1-\delta_{\mathrm{ph}}^6U_6^6
=
\bigl(y^6-\delta_{\mathrm{ph}}^6\bigr)\mathbf 1.
\]
Dividing by $y^6-\delta_{\mathrm{ph}}^6$ gives \Cref{eq:finite-hexagonal-resolvent}.
\end{proof}

Equivalently, one may read \Cref{eq:yukawa-kernel} in the Fourier basis of the six-site ring: $U_6$ is the diagonal form of the step operator, $D_y$ is the shifted walk denominator, and $\mathcal Y_y^{(\varepsilon)}$ is the corresponding regularized Green function. In this reading the admissible path lengths $L_{f,j}$ are graph-distance data on $C_6$, while the residual factors $\Lambda_{f,j}$ encode bounded local transport amplitudes rather than unrestricted fit parameters.

Because $U_6$ has eigenphases $e^{in\pi/3}$, the regularized transport spectrum is
\begin{equation}
\label{eq:transport-spectrum}
\lambda_{y,n}^{-1}
=
y^2+\delta_{\mathrm{ph}}^2-2y\delta_{\mathrm{ph}}\cos\!\left(\frac{n\pi}{3}\right)+\varepsilon^2,
\qquad
n=0,\dots,5,
\end{equation}
for the eigenvalues $\lambda_{y,n}$ of $\mathcal Y_y^{(\varepsilon)}$. In particular,
\begin{lemma}[Finite admissible hexagon gap]
\label{lem:finite-admissible-hexagon-gap}
On the admissible cusp set
\[
y\in\left\{1,\frac13,\frac23\right\},
\]
the smallest singular value of the unregulated hexagon denominator is attained at $y=\frac23$ and $n=0$. Equivalently,
\[
\epsilon_f
:=
\min_{\substack{y\in\{1,\frac13,\frac23\}\\ n=0,\dots,5}}
\left|y-\delta_{\mathrm{ph}}e^{in\pi/3}\right|
=
\left|\frac23-\delta_{\mathrm{ph}}\right|.
\]
Hence the finite admissible hexagon has a strictly positive spectral gap $\epsilon_f$.
\end{lemma}

\begin{proof}
For $\varepsilon=0$, the eigenvalues in \Cref{eq:transport-spectrum} reduce to $|y-\delta_{\mathrm{ph}}e^{in\pi/3}|^{-2}$. On the retained branch one has $\delta_{\mathrm{ph}}\approx \num{0.599882}$, so among the admissible cusps the closest point to the positive real phase is $y=2/3$ with $n=0$. This gives the stated minimum.
\end{proof}

\begin{definition}[Gap-stable admissibility]
\label{def:gap-stable-admissibility}
Let
\[
D_0:=D_{2/3}.
\]
We say that the admissible transport sector is gap-stable if the physical transport operator has the form
\[
D_f=D_0+V_{\mathrm{adm}},
\qquad
[\Padm,V_{\mathrm{adm}}]=0,
\qquad
\|V_{\mathrm{adm}}\|<\frac{\epsilon_f}{2}.
\]
\end{definition}

\begin{proposition}[Lift to the full admissibility-projected operator]
\label{prop:gap-lift-full-operator}
If admissibility is gap-stable in the sense of \Cref{def:gap-stable-admissibility}, then for every
\[
\psi\in\mathrm{Ran}(\Padm)
\]
one has
\[
\|D_f\psi\|
\ge
\frac{\epsilon_f}{2}\|\psi\|.
\]
Consequently
\begin{equation}
\label{eq:gap-lift-full-operator}
\Padm D_f^\dagger D_f \Padm
\ge
\frac{\epsilon_f^2}{4}\,\Padm
>
0.
\end{equation}
\end{proposition}

\begin{proof}
Since $\psi\in\mathrm{Ran}(\Padm)$ and $[\Padm,V_{\mathrm{adm}}]=0$,
\[
\|D_f\psi\|
\ge
\|D_0\psi\|-\|V_{\mathrm{adm}}\psi\|
\ge
\bigl(\epsilon_f-\|V_{\mathrm{adm}}\|\bigr)\|\psi\|
\ge
\frac{\epsilon_f}{2}\|\psi\|.
\]
Squaring gives \Cref{eq:gap-lift-full-operator}.
\end{proof}

\begin{equation}
\label{eq:mobius-quotient}
\frac{\lambda_{y,0}}{\lambda_{y,3}}
=
\frac{(y+\delta_{\mathrm{ph}})^2+\varepsilon^2}{(y-\delta_{\mathrm{ph}})^2+\varepsilon^2}.
\end{equation}
\[
\lim_{\varepsilon\to 0}\frac{\lambda_{y,0}}{\lambda_{y,3}}
=
\left(\frac{y+\delta_{\mathrm{ph}}}{y-\delta_{\mathrm{ph}}}\right)^2
\]
\begin{remark}[Near-pole hierarchy mechanism]
On the retained canonical transport saddle from \Cref{eq:transport-bridge-offset},
\[
\delta_{\mathrm{ph}}\approx \num{0.599882}.
\]
Hence the Möbius quotient contains a built-in near-pole whenever $y$ approaches $\delta_{\mathrm{ph}}$, making large hierarchy ratios structurally plausible. In the closed transport grammar this is the hierarchy mechanism, while the bounded holonomy factors keep it from degenerating into an arbitrary tuning surface.
\end{remark}

\begin{proposition}[Finite positive transport algebra]
\label{prop:finite-positive-transport-algebra}
Let
\[
C:=\frac12\bigl(U_6+U_6^\dagger\bigr),
\qquad
a:=y^2+\delta_{\mathrm{ph}}^2+\varepsilon^2,
\qquad
b:=y\delta_{\mathrm{ph}}.
\]
Then
\[
D_y^\dagger D_y+\varepsilon^2=a\,\mathbf 1-2b\,C.
\]
The cosine operator $C$ has spectrum
\[
\mathrm{Spec}(C)=\left\{1,\frac12,-\frac12,-1\right\}
\]
and satisfies
\begin{equation}
\label{eq:c6-cosine-polynomial}
4C^4-5C^2+\mathbf 1=0.
\end{equation}
Consequently the regularized positive transport kernel reduces to the cubic polynomial
\begin{equation}
\label{eq:finite-positive-transport-kernel}
\mathcal Y_y^{(\varepsilon)}
=
\frac{
a^3-5ab^2
+2b(a^2-5b^2)C
+4ab^2C^2
+8b^3C^3
}{
(a^2-b^2)(a^2-4b^2)
}.
\end{equation}
\end{proposition}

\begin{proof}
Because $U_6$ is unitary,
\[
D_y^\dagger D_y
=
\bigl(y\,\mathbf 1-\delta_{\mathrm{ph}}U_6^\dagger\bigr)
\bigl(y\,\mathbf 1-\delta_{\mathrm{ph}}U_6\bigr)
=
\bigl(y^2+\delta_{\mathrm{ph}}^2\bigr)\mathbf 1
-y\delta_{\mathrm{ph}}(U_6+U_6^\dagger),
\]
which gives the first identity after adding $\varepsilon^2$. The eigenvalues of $C$ are $\cos(n\pi/3)$ for $n=0,\dots,5$, namely $1,\frac12,-\frac12,-1$. Hence $C$ obeys the minimal polynomial \Cref{eq:c6-cosine-polynomial}. The inverse of $a-2bc$ on this four-point spectrum is therefore represented by a unique cubic interpolation polynomial in $c$, and evaluating that interpolation yields \Cref{eq:finite-positive-transport-kernel}.
\end{proof}

\begin{remark}[Carrier rank echo in the transport polynomial]\label{rem:transport-rank-echo}
On the minimal carrier branch $\gcar=5$, the hexagonal cosine polynomial of \Cref{eq:c6-cosine-polynomial} can be rewritten as
\[
(\gcar-1)C^4-\gcar C^2+\mathbf 1=0.
\]
Thus the transport closure polynomial carries the same pair $(\gcar-1,\gcar)$ that governs the hard carrier packet. Equivalently,
\[
4C^4-5C^2+\mathbf 1=(4C^2-\mathbf 1)(C^2-\mathbf 1)=(2C-\mathbf 1)(2C+\mathbf 1)(C-\mathbf 1)(C+\mathbf 1).
\]
So the distinguished value $C=1/2$ reappears on the transport side, mirroring the positive carrier root $Y=1/2$ of the factorized carrier polynomial in Section~3.3. At present this should be read as an exact compression identity on the retained hexagonal branch, not as a derivation of the transport sector from the carrier theorem.
\end{remark}

Absolute masses are attached to admissible path lengths,
\begin{equation}
\label{eq:path-length-yukawa}
y_{f,j}^{\mathrm{UV}}=\lambdaY^{L_{f,j}}\Lambda_{f,j},
\qquad
\lambdaY=\sqrt{\phiz(1-\phiz)}.
\end{equation}
The stronger matrix-level statement, used as the main transport holonomy grammar below, is that
\begin{equation}
\label{eq:yukawa-holonomy-matrix}
(Y_f)_{ij}
=
\lambdaY^{L^L_{f,i}+L^R_{f,j}}
\left\langle e_i,\mathrm{Hol}_{\Gamma_{ij}}(\nabla_F)e_j\right\rangle.
\end{equation}
\begin{proposition}[Holonomy triangle phase]
\label{prop:holonomy-triangle-phase}
Define the minimal closed family holonomy by
\[
\Pi_\triangle
:=
\mathrm{Hol}_{\Gamma_{12}}(\nabla_F)\,
\mathrm{Hol}_{\Gamma_{23}}(\nabla_F)\,
\mathrm{Hol}_{\Gamma_{31}}(\nabla_F),
\qquad
\delta_{\mathrm{CKM}}
:=
\operatorname{Arg}\bigl(\tr\Pi_\triangle\bigr).
\]
On the retained transport branch the discrete base holonomy satisfies
\[
\operatorname{Arg}\bigl(\tr\Pi_\triangle^{(0)}\bigr)=\frac{\pi}{3},
\]
and the smallest admissible nonabelian triangle-area contribution is
\[
A_\triangle=3\lambda_C^2+O(\lambda_C^4).
\]
Therefore
\begin{equation}
\label{eq:ckm-holonomy-phase}
\delta_{\mathrm{CKM}}
=
\frac{\pi}{3}+3\lambda_C^2+O(\lambda_C^4).
\end{equation}
\end{proposition}

\begin{proof}[Proof sketch]
The family holonomy is already encoded in \Cref{eq:yukawa-holonomy-matrix}. Closing the minimal triangle isolates the gauge-invariant phase information carried by the family connection. The $\mathbb Z_6$ transport grammar fixes the base phase at $\pi/3$, while the first nontrivial nonabelian area correction is quadratic in the Cabibbo seed. Keeping only the leading admissible area term gives \Cref{eq:ckm-holonomy-phase}.
\end{proof}

\Cref{eq:path-length-yukawa} is the diagonal or singular-value reduction of \Cref{eq:yukawa-holonomy-matrix}, not the conceptual starting point of the transport sector.
After the geometric electroweak closure fixes $\vgeo$, the intrinsic compiler mass rows are emitted through
\begin{equation}
\label{eq:mass-bridge}
\hat m_{f,j}:=\frac{y_{f,j}\vgeo}{\sqrt2}.
\end{equation}
These hatted quantities are the intrinsic compiler masses on the retained branch. The observation-surface masses are obtained only after the declared comparison map:
\[
m_{f,j}^{\mathrm{obs}}:=\mathcal R_{\SM}(\hat m_{f,j}).
\]
The role of the residual factors $\Lambda_{f,j}$ must be kept explicit. They are bounded $O(1)$ transport prefactors, not unrestricted fit knobs. The RG bookkeeping behind direct $\vgeo$ insertions is part of the declared comparison map discussed later in \Cref{sec:rg-map-compiler-outputs}.

\begin{remark}[Singlet-projected cusp lift to the transport hexagon]
The transport cusps do not come from the raw two-point carrier spectrum alone. The carrier itself has
\[
\mathrm{Spec}(Y|_E)=\left\{-\frac13,\frac12\right\}.
\]
After exterior lift to the one-family packet,
\[
\mathrm{Spec}(Y|_{S^+})
=
\left\{0,-\frac23,\frac16,1,-\frac12,\frac13\right\},
\]
with multiplicities displayed in \Cref{tab:one-family}. Restricting to the singlet channels and taking absolute values yields
\[
\{|Y(e^c)|,\ |Y(u^c)|,\ |Y(d^c)|\}
=
\left\{1,\frac23,\frac13\right\},
\]
which are exactly the admissible cusp values of the transport hexagon. In this sense the transport sector reads the singlet-projected, family-lifted hypercharge spectrum through the resolvent
\[
D_y=y\mathbf 1-\delta_{ph}U_6
\]
on $C_6$. The transport grammar therefore lives on the same rational grid as the carrier packet, but only after exterior lift and singlet projection.
\end{remark}
